\section{Context of Electron-Positron Alliance}
\label{sec:context}

Modern collider datasets provide unprecedented statistical power, and experimental uncertainties on many Standard Model measurements are now at the percent level. In this regime, theoretical uncertainties increasingly limit the overall precision, and those associated with Quantum Chromodynamics (QCD) are among the dominant contributions. To achieve the physics goals of future programs, including the High-Luminosity LHC (HL-LHC) and Future Circular Collider (FCC-ee), theoretical predictions must therefore reach a precision comparable to the expected experimental uncertainties, in many cases at the sub-percent level, so that potential new physics anomalies can be interpreted reliably.

Precision QCD studies can be pursued at hadron colliders such as the LHC, but they are most cleanly and systematically tested in lepton collider environments. At present, however, there is no active high-energy lepton collider operating at Large Electron--Positron collider (LEP) energies or above. To address this, the Electron-Positron Alliance (EPA)~\cite{EEAlliance:website} works with LEP collaboration members to curated and reanalyze archived datasets, opening new opportunities enabled by three decades of theoretical advances and modern experimental techniques. The combination of ultra-precise $e^+e^-$ data and state-of-the-art perturbative and non-perturbative theory tools provides a uniquely powerful environment for high-precision studies. The EPA has produced a series of precision measurements using these archived datasets, including two-particle correlations~\cite{Badea:2019vey}, jet energy spectra and substructure~\cite{Chen:2021uws}, long-range near-side correlations at $WW$ threshold~\cite{Chen:2024lrc}, energy-energy correlators~\cite{Bossi:2025eec}, and unbinned measurement of thrust~\cite{Badea:2025ubt}.

More broadly, this program is a compelling opportunity for AI for Science in a domain where precision is paramount. While AI is increasingly used to assist theoretical calculations and generate predictions, in the physical sciences such predictions ultimately require experimental validation. In high-energy physics, this validation depends on collider datasets that are rarely publicly accessible. The $e^+e^-$ environment, particularly at LEP energies, provides the cleanest setting for precision tests of the Standard Model (SM), with well-defined initial states and calculable dynamics. The successful curation and use of archived open data therefore establishes the conditions for a theory--experiment loop in which AI systems, guided by domain experts, can propose new calculations, observables, or improved models and directly confront them with high-precision data. This creates a concrete framework for developing AI systems that do not merely interpolate existing knowledge, but actively refine theoretical descriptions through direct comparison with experiment. We believe that this technology is the start of an exciting scientific revolution.

\subsection*{Lessons from AI--Physicist Collaboration}

Carrying out this analysis has also produced practical insights into how AI agents and physicists can work together effectively. Our workflow followed a constrained human-in-the-loop cycle in which the physicist defined physics objectives, acceptance criteria, and interpretation authority, while the AI agent owned implementation and execution. Crucially, experienced physicists contribute domain context that cannot yet be reduced to formal specifications alone, including judgments about physically meaningful binning, acceptable regularization behavior, uncertainty communication, and when apparent anomalies warranted deeper investigation. A major future goal identified in this work is to translate this expert knowledge into explicit, code-based verification layers. 
%Requests were made concrete and executable, each iteration linked to specific scripts, outputs, and note sections, and corrections were applied incrementally rather than batched into large changes.

From this experience, several recurring challenge classes emerged. Implicit physics priors, expert expectations about acceptable binning, regularization behavior, or uncertainty presentation, were often assumed rather than explicitly stated, requiring iterative clarification. Figure formatting and style choices (label offsets, legend placement, pad geometry) required many short feedback loops, with cosmetic shortcomings sometimes affecting interpretability and therefore physics conclusions. Synchronization across rapidly evolving scripts, outputs, and note text demanded continuous provenance refreshing. These lessons motivate creation of reusable tools, stronger configuration contracts, clearer channels for encoding expert heuristics into machine readable constraints in future iterations of this program, and new verification layers throughout the autonomous analysis flow.